\subsection{Year Fraction Convention}\label{sec:yearfrac}

Interest income and other calculations are based on \textit{per annum} interest rates. Therefore, the year-fraction function interface $Y$ is used to calculate the \textit{fraction of a year} between any two times $s$ and $t$ with $t>s$ for which e.g. an (per annum) interest rate applies according to some day count convention \attr{DCC}

\[
	\yfrfunc: s,t,\attr{DCC} \mapsto \Real
\]

Note, the year fraction function interface only defines the structure of year fraction functions but not an actual implementation thereof, or the respective \attr{DCC}, respectively. Therefore, any \attr{DCC} can be implemented according to the interface above supporting user-defined year fraction functions.\\

For brevity we will omit the \attr{DCC} function argument wherever this does not lead to confusion.


\subsection{Contract Role Sign Convention}\label{sec:cntrl}

The two parties to a contract are defined through attributes \attr{CRID} and \attr{CPID}. The first is the party initially \textit{creating} the contract and the second is the counterparty, respectively. Thereby, both \attr{CRID}/\attr{CPID} can take any \textit{role} in the contract or, more specifically, they can be the lender or borrower in a loan (PAM), fixed receiver or payer in an interest rate swap (SWAPS), etc.\\
